\documentclass[10pt]{article}

\usepackage{graphicx}
\usepackage[margin=1in]{geometry}
\usepackage{parskip}
\usepackage{xr}

\externaldocument{main}

\begin{document}

\appendix
\section{Label Placement}
\label{sec:label-placement}

While VERA exracts semantically meaningful visual explanations, there remains the practical apsect of how to visualize these visual explanations on-screen. A key challenge in this is how to position region labels, so that each region polygon is accompanied by a textual description of what the region represents.
Label placement has been a well-studies problem, with several approaches being proposed over the years. However, these algorithms typically deal with annotating specific points on a scatter plot and are not directly applicable to our task here.
However, extending any one of these is straightforward to our particular task.

\noindent \textbf{Label Placement Algorithm.} We here describe a force-directed label placement algorithm we designed for VERA. A good label positioning for a VERA explanation must necessarily fulfill the following conditions:
\begin{itemize}
  \item Labels must be inside the axes bounding box and must not overflow.
  \item Labels must not overlap one another as overlapping text is often completely illegible.
  \item Labels should not overlap with the regions, since labels may be illegible.
  \item Labels should be positioned closely to their associated region.
\end{itemize}

\noindent \textbf{Drawbacks and Limitations.} We here describe the algorithm.

\end{document}
